\documentclass{article}
\usepackage[T1]{fontenc}
\usepackage{lmodern}
\usepackage{graphicx}
\usepackage{amsmath,amssymb}
\usepackage[a4paper,margin=2cm]{geometry}
\usepackage[absolute,overlay]{textpos}
\setlength{\TPHorizModule}{1mm}
\setlength{\TPVertModule}{1mm}
\usepackage[indonesian]{babel}
\usepackage{hyperref}
\setlength{\emergencystretch}{2em}
% unit: Halaman 4

\begin{document}

% === Halaman 4 (python/native) ===

4

Teori Informasi (Information Theory)

• Information theory: cabang ilmu yang mempelajari kuantisasi, 

penyimpanan, transmisi, dan pengolahan informasi

• Contoh kuantisasi: 

 - 1 bit untuk mengkodekan jenis kelamin (M/F)

 - 3 bit untuk mengkodekan nama hari  (ada 7 hari)

 - 4 bit untuk mengkodekan  0 s/d 9 

• Contoh penyimpanan: pemampatan data untuk mengurangi ukuran ruang storage

\end{document}
